% ============================ CHAPITRE 1 ==============================
\chapter{Cadre général du projet}

\section*{Introduction}
Ce premier chapitre situe le projet dans son contexte. Il délimite d'abord le
cadre du projet, présente ensuite l'organisme d'accueil et ses missions, puis
analyse la situation existante afin d'en dégager la problématique et la solution
proposée. Il se termine par la présentation de la \textbf{démarche} et de la
\textbf{planification} du projet.

% ----------------------------------------------------------------------
\section{Cadre du projet}
Le présent travail constitue un \textbf{projet de fin d'études} réalisé en vue de
l'obtention de la Licence en Business Computing, parcours Business Intelligence, à
l'\textbf{Esprit School of Business}. Il s'est déroulé dans le cadre d'un stage au
sein de la \textbf{BTK -- Banque Tuniso-Koweïtienne}, du \textbf{16 février au
16 juin 2026}. Le projet porte sur la conception et le développement d'une
application web de gestion du réseau d'agences bancaires, adossée à un volet
décisionnel d'aide à la décision.

Les objectifs assignés à ce projet sont les suivants :
\begin{itemize}[noitemsep]
  \item \textbf{centraliser} la gestion des agences, des employés, des clients et
        des objectifs commerciaux dans une base unique ;
  \item \textbf{sécuriser} l'accès par une gestion fine des rôles et des droits ;
  \item \textbf{formaliser} les opérations sensibles par des circuits de validation
        et en assurer la traçabilité ;
  \item \textbf{suivre la présence} des employés au moyen d'un module de pointage ;
  \item \textbf{restituer} des indicateurs d'aide à la décision via une chaîne ETL,
        un tableau de bord et des rapports Power BI ;
  \item \textbf{valoriser} les données par une segmentation des agences fondée sur
        l'apprentissage automatique.
\end{itemize}

% ----------------------------------------------------------------------
\section{Présentation de l'organisme d'accueil}
La \textbf{Banque Tuniso-Koweïtienne} (BTK Bank) est une banque universelle de
droit tunisien, créée le 25 février 1981 dans le cadre d'une convention de
coopération bilatérale entre la Tunisie et le Koweït. Constituée en société
anonyme et régie par la loi bancaire n\textsuperscript{o}~2016-48, elle dispose
d'un capital social de 200~millions de dinars et a son siège social à Tunis. À
travers un réseau d'agences couvrant l'ensemble du territoire tunisien, elle
accompagne aussi bien les particuliers que les entreprises.

\begin{figure}[H]
  \centering
  
\includegraphics[width=0.18\textwidth]{logo_btk.png}
  \caption{Logo de la BTK -- Banque Tuniso-Koweïtienne}
  \label{fig:logo_btk}
\end{figure}

Le tableau~\ref{tab:fiche_btk} en résume la fiche signalétique.

\begin{table}[htbp]
\centering
\caption{Fiche signalétique de la BTK}
\label{tab:fiche_btk}
\renewcommand{\arraystretch}{1.2}
\begin{tabularx}{\textwidth}{>{\bfseries}p{4.2cm}X}
\toprule
\textmd{\textbf{Élément}} & \textbf{Renseignement}\\
\midrule
Dénomination sociale & Banque Tuniso-Koweïtienne (BTK Bank)\\
Forme juridique & Société Anonyme (S.A.)\\
Date de création & 25 février 1981\\
Capital social & 200\,000\,000 DT\\
Siège social & 10 bis, Avenue Mohamed V, 1001 Tunis\\
Secteur d'activité & Banque universelle\\
Actionnaire de référence & Groupe Elloumi (60 \%)\\
Réseau & 5 directions régionales, 6 succursales, 45 agences\\
Effectif (2024) & $\approx$ 400 collaborateurs\\
\bottomrule
\end{tabularx}
\end{table}

\subsection{Missions}
En tant que banque universelle, la BTK assure un ensemble de missions couvrant
les principaux métiers bancaires :
\begin{itemize}[noitemsep]
  \item \textbf{Banque de détail} : comptes, cartes, crédits à la consommation et
        immobiliers, épargne et banque digitale destinés aux particuliers ;
  \item \textbf{Banque des entreprises} : financement de l'investissement et du
        fonctionnement, financement du commerce extérieur et gestion de
        trésorerie ;
  \item \textbf{Opérations internationales et de marché} : change, couverture du
        risque de change et placements ;
  \item \textbf{Innovation et digitalisation} : modernisation des services et
        offre en ligne.
\end{itemize}

\subsection{Objectifs}
À travers ces missions, la BTK poursuit plusieurs objectifs stratégiques :
\begin{itemize}[noitemsep]
  \item soutenir le tissu économique en finançant les entreprises et les projets
        d'investissement ;
  \item proposer des solutions financières adaptées aux ménages ;
  \item contribuer à la création d'emplois et à la dynamique de croissance
        nationale ;
  \item consolider sa solidité financière et accélérer sa transformation
        digitale.
\end{itemize}

% ----------------------------------------------------------------------
\section{Étude de l'existant}

\subsection{Description de la situation existante}
La gestion actuelle du réseau repose en grande partie sur des traitements manuels
et des fichiers bureautiques non centralisés. Les informations relatives aux
agences, aux employés et aux clients sont saisies et conservées de manière
hétérogène ; les demandes de création ou de modification circulent par des canaux
informels ; et la présence des employés n'est suivie par aucun dispositif
formalisé.

\subsection{Critique de la situation existante}
Cette organisation présente plusieurs limites :
\begin{itemize}[noitemsep]
  \item dispersion et redondance des données, sources d'incohérences ;
  \item absence de circuit de validation formalisé pour les opérations sensibles ;
  \item absence de suivi de présence (pointage) des employés ;
  \item faible traçabilité des actions et des décisions ;
  \item difficulté à produire des indicateurs consolidés pour le pilotage ;
  \item gestion des droits d'accès insuffisamment maîtrisée.
\end{itemize}

\subsection{Problématique}
La principale problématique qui se pose est la suivante : comment
\textbf{centraliser et fiabiliser} la gestion du réseau d'agences bancaires,
\textbf{encadrer} les opérations sensibles par des circuits de validation,
\textbf{assurer le suivi de présence} des employés, et \textbf{exploiter} les
données ainsi structurées pour fournir aux responsables des indicateurs d'aide à
la décision ?

\subsection{Solution proposée}
La solution retenue consiste en une \textbf{application web centralisée},
structurée autour d'une base de données Oracle, qui couvre la gestion des agences,
des employés, des clients et des objectifs, gère les rôles et la sécurité des
accès, assure le \textbf{pointage} de présence des employés, et formalise les
demandes de création et de modification au moyen de \textbf{workflows de
validation}. Un \textbf{volet décisionnel} complète le dispositif : une chaîne ETL
consolide les données dans un datamart qui alimente à la fois Microsoft Power BI,
un tableau de bord embarqué et une \textbf{segmentation des agences} par
apprentissage automatique.

\subsection{Justification du choix de la solution}
Le choix d'une application web centralisée se justifie par plusieurs facteurs :
l'\textbf{accessibilité} (un simple navigateur suffit, sans installation sur les
postes), la \textbf{centralisation} des données dans une base unique garantissant
leur cohérence, la \textbf{sécurité} offerte par une gestion fine des rôles et des
workflows de validation, et l'\textbf{ouverture décisionnelle} permise par
l'exposition des données à des outils d'analyse et de \emph{reporting}.

% ----------------------------------------------------------------------
\section{Méthodologie et planification du projet}

\subsection{Démarche de développement}
Le périmètre du projet ayant été précisé progressivement au fil des échanges avec
l'organisme d'accueil, une \textbf{démarche itérative et incrémentale}, d'inspiration
\emph{agile}, a été adoptée. Plutôt qu'un enchaînement séquentiel figé, cette
approche privilégie des livraisons fonctionnelles successives et une validation
régulière avec l'encadrement, ce qui permet d'\textbf{adapter la solution} aux
retours et de \textbf{détecter les risques au plus tôt}. Le travail a été organisé
en phases classiques du génie logiciel --- analyse et spécification des besoins,
conception, réalisation, puis volet décisionnel --- reprises comme fil conducteur
des chapitres de ce rapport.

\subsection{Planification}
Le projet s'est déroulé sur quatre mois, du 16 février au 16 juin 2026. La
figure~\ref{fig:gantt} en présente le planning prévisionnel sous forme de
diagramme de Gantt, organisé selon les grandes phases du projet.

\begin{figure}[H]
  \centering
  \includegraphics[width=\textwidth]{gantt.png}
  \caption{Planning prévisionnel du projet (diagramme de Gantt)}
  \label{fig:gantt}
\end{figure}

\section*{Conclusion}
Ce chapitre a posé le cadre du projet : l'organisme d'accueil et ses missions,
l'étude de l'existant qui a mis en évidence la problématique, la solution
proposée, ainsi que la démarche et la planification retenues. Le chapitre suivant
est consacré à l'\textbf{analyse et à la spécification des besoins}.
